El objetivo general de este trabajo es construir el \dfn{frontend} del sistema SCDEE,
integrado con el contrato real de su \dfn{backend}, y documentar el proceso de desarrollo
asistido por agentes de programación de forma que pueda generalizarse. Se concreta en los
objetivos específicos siguientes, cada uno verificable con la evidencia que se indica.

\begin{description}
  \item[O1. Interfaz para todos los actores.] Construir una aplicación web de página única
        (\acs{spa}) que cubra los flujos de los cuatro actores del sistema: gestión de la
        organización, coordinación de asignatura, docencia y estudiantado.
        \emph{Evidencia:} \cref{SEC:AN-ACTORS,CAP:DESIGN,CAP:IMPLEMENTATION}.
  \item[O2. Fidelidad al contrato.] Integrar la interfaz con la especificación OpenAPI del
        \dfn{backend} sin inventar operaciones ni campos, y documentar como propuestas los
        datos que la interfaz necesita y el contrato no ofrece.
        \emph{Evidencia:} \cref{SEC:AN-BACKEND,SEC:AN-GAPS,SEC:TEST-COVERAGE,AP:API}.
  \item[O3. Desarrollo independiente del servidor.] Disponer de una \dfn{mock} que responda
        con los mismos formatos y reglas de negocio que el servidor real, de modo que el
        mismo código funcione contra ambos.
        \emph{Evidencia:} \cref{SEC:DES-DATA,SEC:TEST-REAL}.
  \item[O4. Calidad verificable.] Acompañar el código de pruebas de flujo automatizadas y
        de comprobaciones deterministas de estilo y de reparto del trabajo que se ejecuten
        antes de cada \textit{commit}.
        \emph{Evidencia:} \cref{CAP:TESTING}.
  \item[O5. Lenguaje visual educativo y accesible.] Definir un sistema de diseño basado en
        \textit{tokens} y componentes primitivos que garantice una apariencia coherente y
        legible, tomando como referencia las pautas \acs{wcag} \cite{wcag22}.
        \emph{Evidencia:} \cref{SEC:DES-VISUAL,AP:STYLE}.
  \item[O6. Análisis del proceso agéntico.] Registrar el proceso (decisiones, errores y
        correcciones), declarar con transparencia el uso de herramientas de \acs{ia} y
        analizar qué partes de la construcción de un \dfn{frontend} a partir de un
        \dfn{backend} son deterministas, cuáles requieren razonamiento y cuáles requieren
        a una persona.
        \emph{Evidencia:} \cref{SEC:IMP-AI,CAP:GENERALIZATION,AP:ITERATIONS,AP:AI}.
\end{description}

\pendiente{Validar la lista de objetivos con el tutor.}
